%% Chapter 9 — Mathematical Modelling
\chapter{Mathematical Modelling}
\label{ch:modelling}

\section{Introduction}

This chapter documents the mathematical and computational models underpinning the \BDE{} pipeline, including the Python implementation, uncertainty quantification framework, and spatial autocorrelation analysis. Full source code is maintained in the DreamLab AI repository; key algorithms are reproduced here for transparency and reproducibility.


\section{Software Environment}

\begin{table}[H]
\centering
\caption{Software dependencies and versions.}
\label{tab:software-env}
\begin{tabular}{llp{60mm}}
\toprule
\textbf{Package} & \textbf{Version} & \textbf{Purpose} \\
\midrule
Python              & 3.11.8    & Runtime environment \\
scikit-learn        & \dataplaceholder{VER} & Random Forest classifier, cross-validation \\
NumPy               & \dataplaceholder{VER} & Numerical computation \\
pandas              & \dataplaceholder{VER} & Tabular data manipulation \\
GeoPandas           & \dataplaceholder{VER} & Spatial data handling \\
rasterio            & \dataplaceholder{VER} & Raster I/O \\
rasterstats         & \dataplaceholder{VER} & Zonal statistics \\
PySAL (libpysal)    & \dataplaceholder{VER} & Spatial autocorrelation analysis \\
SciPy               & \dataplaceholder{VER} & Statistical distributions, integration \\
Matplotlib          & \dataplaceholder{VER} & Visualisation \\
Google Earth Engine  & \dataplaceholder{VER} & Remote sensing data access \\
GDAL                & \dataplaceholder{VER} & Geospatial data processing \\
\bottomrule
\end{tabular}
\end{table}


\section{Classification Model}

\subsection{Feature Matrix Construction}

For each parcel $i$ in the study area, a feature vector $\mathbf{x}_i \in \mathbb{R}^p$ is constructed from the following components:

\begin{equation}
\mathbf{x}_i = \left[ \overline{\text{NDVI}}_i,\; \sigma_{\text{NDVI},i},\; \overline{\text{EVI}}_i,\; \overline{\text{NDMI}}_i,\; \overline{\text{SAVI}}_i,\; \mathbf{t}_i,\; \mathbf{g}_i,\; \mathbf{e}_i,\; \mathbf{m}_i \right]
\label{eq:feature-vector}
\end{equation}

\noindent where:
\begin{description}[nosep]
  \item[$\mathbf{t}_i$] Temporal features: NDVI amplitude, peak date, greenup rate, senescence rate;
  \item[$\mathbf{g}_i$] GLCM texture features: contrast, dissimilarity, homogeneity, ASM, entropy, correlation;
  \item[$\mathbf{e}_i$] Topographic features: elevation, slope, aspect, topographic wetness index;
  \item[$\mathbf{m}_i$] MasterMap features: parcel area, perimeter-to-area ratio, OS descriptive group encoding.
\end{description}

The total feature dimensionality is $p = \dataplaceholder{P}$.

\subsection{Random Forest Implementation}

\begin{lstlisting}[language=Python,caption={Random Forest classifier training pipeline.},label={lst:rf-train}]
import numpy as np
from sklearn.ensemble import RandomForestClassifier
from sklearn.model_selection import GroupKFold
from sklearn.metrics import classification_report

def train_habitat_classifier(X_train, y_train, groups, params):
    """Train RF classifier with spatial cross-validation."""
    clf = RandomForestClassifier(
        n_estimators=params['n_estimators'],
        max_depth=params['max_depth'],
        min_samples_leaf=params['min_samples_leaf'],
        max_features='sqrt',
        class_weight='balanced',
        oob_score=True,
        random_state=42,
        n_jobs=-1
    )

    # Spatial cross-validation to prevent autocorrelation leakage
    cv = GroupKFold(n_splits=5)
    cv_scores = []
    for fold, (train_idx, val_idx) in enumerate(
        cv.split(X_train, y_train, groups)
    ):
        clf_fold = clone(clf)
        clf_fold.fit(X_train[train_idx], y_train[train_idx])
        score = clf_fold.score(X_train[val_idx], y_train[val_idx])
        cv_scores.append(score)

    # Final model on full training set
    clf.fit(X_train, y_train)
    return clf, cv_scores
\end{lstlisting}

\subsection{Feature Importance}

\begin{figure}[H]
\centering
\begin{tikzpicture}
  \draw[dreamlab-light,fill=dreamlab-light,rounded corners] (0,0) rectangle (12,7);
  \node[dreamlab-grey,font=\sffamily\large] at (6,3.5) {Figure placeholder: Feature importance bar chart};
  \node[dreamlab-grey,font=\sffamily\footnotesize] at (6,2.5) {../output/figures/feature\_importance.png};
\end{tikzpicture}
\caption{Random Forest feature importance (Gini impurity decrease) for the top 20 features. Spectral indices and temporal features dominate, with NDVI mean, EVI mean, and NDVI amplitude as the three most informative features.}
\label{fig:feature-importance}
\end{figure}


\section{BM\,4.0 Scoring Engine}

The scoring engine implements \cref{eq:bu} programmatically, applying the statutory lookup tables for distinctiveness, condition multipliers, and strategic significance.

\begin{lstlisting}[language=Python,caption={BM\,4.0 biodiversity unit calculation.},label={lst:bu-calc}]
import pandas as pd

# Distinctiveness lookup (UKHab Level 3 -> score)
DISTINCTIVENESS = {
    'g4': 2, 'g3c': 4, 'g1a': 6, 'g3a': 8,
    'w1f': 6, 'w1g': 4, 'w2c': 4,
    'h1a': 6, 'c1': 2, 'c1a': 4,
    'r1': 4, 'u1b': 0, 'u1c': 0, 'u1e': 0,
    'g6': 2,
}

# Condition multiplier lookup
CONDITION_MULTIPLIER = {
    'Good': 3.0, 'Fairly Good': 2.5, 'Moderate': 2.0,
    'Fairly Poor': 1.5, 'Poor': 1.0, 'N/A': 0.0,
}

# Strategic significance multiplier
STRATEGIC_MULTIPLIER = {
    'High': 1.15, 'Medium': 1.10, 'Low': 1.00,
}

def calculate_bu(parcels: pd.DataFrame) -> pd.Series:
    """Calculate biodiversity units for each parcel."""
    bu = (
        parcels['area_ha']
        * parcels['ukhab_class'].map(DISTINCTIVENESS)
        * parcels['condition'].map(CONDITION_MULTIPLIER)
        * parcels['strategic'].map(STRATEGIC_MULTIPLIER)
    )
    return bu
\end{lstlisting}


\section{Monte Carlo Uncertainty Framework}

\subsection{Perturbation Model}

The Monte Carlo simulation perturbs three parameters simultaneously at each iteration $m \in \{1, \ldots, M\}$ where $M = 10{,}000$:

\begin{align}
\tilde{y}_i^{(m)} &\sim \text{Categorical}\!\left(\hat{\mathbf{p}}_i\right) \label{eq:mc-class} \\
\tilde{c}_i^{(m)} &= c_i + \epsilon_{c}^{(m)}, \quad \epsilon_{c}^{(m)} \sim \mathcal{N}(0, \sigma_c^2) \label{eq:mc-cond} \\
\tilde{A}_i^{(m)} &= A_i \cdot (1 + \epsilon_{A}^{(m)}), \quad \epsilon_{A}^{(m)} \sim \mathcal{U}(-0.02, 0.02) \label{eq:mc-area}
\end{align}

\noindent where $\hat{\mathbf{p}}_i$ is the Random Forest's predicted class probability vector for parcel $i$, $\sigma_c$ is the standard deviation of condition calibration residuals (\dataplaceholder{VALUE}), and the area perturbation is uniform $\pm$2\%.

\subsection{Implementation}

\begin{lstlisting}[language=Python,caption={Monte Carlo delta simulation.},label={lst:mc-sim}]
import numpy as np
from typing import NamedTuple

class MCResult(NamedTuple):
    deltas: np.ndarray  # shape (M,)
    ci_low: float
    ci_high: float
    p_gain: float

def monte_carlo_delta(
    parcels_t0, parcels_t1, class_probs_t0, class_probs_t1,
    sigma_c: float, M: int = 10_000, seed: int = 42
) -> MCResult:
    """Run Monte Carlo simulation for total delta BU."""
    rng = np.random.default_rng(seed)
    n = len(parcels_t0)
    deltas = np.empty(M)

    for m in range(M):
        # Perturb classifications
        classes_t0 = np.array([
            rng.choice(len(p), p=p) for p in class_probs_t0
        ])
        classes_t1 = np.array([
            rng.choice(len(p), p=p) for p in class_probs_t1
        ])

        # Perturb condition (clipped to valid range)
        cond_noise = rng.normal(0, sigma_c, size=n)

        # Perturb area
        area_noise = rng.uniform(-0.02, 0.02, size=n)

        # Calculate perturbed BU at both epochs
        bu_t0 = calculate_perturbed_bu(
            parcels_t0, classes_t0, cond_noise, area_noise
        )
        bu_t1 = calculate_perturbed_bu(
            parcels_t1, classes_t1, cond_noise, area_noise
        )

        deltas[m] = bu_t1.sum() - bu_t0.sum()

    return MCResult(
        deltas=deltas,
        ci_low=np.percentile(deltas, 2.5),
        ci_high=np.percentile(deltas, 97.5),
        p_gain=np.mean(deltas > 0)
    )
\end{lstlisting}

\subsection{Convergence Analysis}

Convergence of the Monte Carlo estimator was assessed by monitoring the running mean and 95\% CI width as a function of iteration count:

\begin{equation}
\hat{\mu}_m = \frac{1}{m} \sum_{j=1}^{m} \Delta\text{BU}^{(j)}, \quad \hat{w}_m = \hat{q}_{0.975,m} - \hat{q}_{0.025,m}
\label{eq:mc-convergence}
\end{equation}

Convergence was declared when $|\hat{w}_m - \hat{w}_{m-100}| / \hat{w}_{m-100} < 0.001$ for 100 consecutive iterations.


\section{Spatial Autocorrelation Analysis}

\subsection{Motivation}

Spatial autocorrelation in the delta values ($\Delta\text{BU}_i$) would indicate that gains or losses are spatially clustered rather than randomly distributed. Understanding this spatial structure is important for:
\begin{itemize}[nosep]
  \item Identifying landscape-scale drivers of change;
  \item Assessing the reliability of parcel-level independence assumptions;
  \item Informing spatial targeting of future habitat enhancement interventions.
\end{itemize}

\subsection{Global Moran's I}

The global Moran's I statistic tests for overall spatial autocorrelation:

\begin{equation}
I = \frac{N}{\sum_{i}\sum_{j} w_{ij}} \cdot \frac{\sum_{i}\sum_{j} w_{ij}(z_i - \bar{z})(z_j - \bar{z})}{\sum_{i}(z_i - \bar{z})^2}
\label{eq:morans-i}
\end{equation}

\noindent where $z_i = \Delta\text{BU}_i$, $\bar{z}$ is the mean delta, $w_{ij}$ is the spatial weight between parcels $i$ and $j$ (Queen contiguity), and $N$ is the number of parcels.

\begin{table}[H]
\centering
\caption{Global Moran's I results for $\Delta$BU spatial autocorrelation.}
\label{tab:morans-i}
\begin{tabular}{lr}
\toprule
\textbf{Statistic} & \textbf{Value} \\
\midrule
Moran's I                      & \dataplaceholder{I} \\
Expected I (under null)        & \dataplaceholder{E} \\
Variance                       & \dataplaceholder{VAR} \\
Z-score                        & \dataplaceholder{Z} \\
p-value (permutation, $N=999$) & \dataplaceholder{P} \\
\bottomrule
\end{tabular}
\end{table}

\begin{confidencebox}{HIGH}
A Moran's I of \dataplaceholder{I} (p $<$ \dataplaceholder{P}) indicates \dataplaceholder{INTERPRETATION — significant positive/no significant} spatial autocorrelation in the delta values, suggesting that \dataplaceholder{INTERPRETATION}.
\end{confidencebox}

\subsection{Local Indicators of Spatial Association (LISA)}

Local Moran's I statistics identify hotspots and coldspots of biodiversity change:

\begin{equation}
I_i = \frac{(z_i - \bar{z})}{\sigma_z^2} \sum_{j} w_{ij}(z_j - \bar{z})
\label{eq:lisa}
\end{equation}

\begin{figure}[H]
\centering
\begin{tikzpicture}
  \draw[dreamlab-light,fill=dreamlab-light,rounded corners] (0,0) rectangle (12,8);
  \node[dreamlab-grey,font=\sffamily\large] at (6,4) {Figure placeholder: LISA cluster map};
  \node[dreamlab-grey,font=\sffamily\footnotesize] at (6,3) {../output/figures/lisa\_cluster\_map.png};
\end{tikzpicture}
\caption{LISA cluster map for $\Delta$BU values. High-High clusters (red) indicate hotspots of biodiversity gain; Low-Low clusters (blue) indicate coldspots of loss; High-Low and Low-High (light shading) indicate spatial outliers.}
\label{fig:lisa-map}
\end{figure}


\section{Spectral Condition Index Calibration}

\subsection{Calibration Methodology}

The SCI weights ($w_1 \ldots w_4$ in \cref{eq:sci}) were calibrated using ordinary least squares regression against field-surveyed condition scores:

\begin{equation}
\hat{c}_i = w_0 + w_1 \overline{\text{NDVI}}_i + w_2 \sigma_{\text{NDVI},i} + w_3 H_{\text{GLCM},i} + w_4 \text{NDVI}_{\text{amp},i} + \epsilon_i
\label{eq:sci-regression}
\end{equation}

\subsection{Calibration Results}

\begin{table}[H]
\centering
\caption{SCI calibration regression coefficients.}
\label{tab:sci-calibration}
\begin{tabular}{lrrrr}
\toprule
\textbf{Feature} & \textbf{Coefficient} & \textbf{Std.~Error} & \textbf{t-stat} & \textbf{p-value} \\
\midrule
Intercept                       & \dataplaceholder{W0} & \dataplaceholder{SE} & \dataplaceholder{T} & \dataplaceholder{P} \\
$\overline{\text{NDVI}}$        & \dataplaceholder{W1} & \dataplaceholder{SE} & \dataplaceholder{T} & \dataplaceholder{P} \\
$\sigma_{\text{NDVI}}$          & \dataplaceholder{W2} & \dataplaceholder{SE} & \dataplaceholder{T} & \dataplaceholder{P} \\
$H_{\text{GLCM}}$               & \dataplaceholder{W3} & \dataplaceholder{SE} & \dataplaceholder{T} & \dataplaceholder{P} \\
$\text{NDVI}_{\text{amp}}$      & \dataplaceholder{W4} & \dataplaceholder{SE} & \dataplaceholder{T} & \dataplaceholder{P} \\
\midrule
\multicolumn{2}{l}{$R^2$ (adjusted)} & \multicolumn{3}{r}{\dataplaceholder{R2}} \\
\multicolumn{2}{l}{RMSE} & \multicolumn{3}{r}{\dataplaceholder{RMSE}} \\
\multicolumn{2}{l}{N (calibration parcels)} & \multicolumn{3}{r}{\dataplaceholder{N}} \\
\bottomrule
\end{tabular}
\end{table}

\begin{figure}[H]
\centering
\begin{tikzpicture}
  \draw[dreamlab-light,fill=dreamlab-light,rounded corners] (0,0) rectangle (12,7);
  \node[dreamlab-grey,font=\sffamily\large] at (6,3.5) {Figure placeholder: SCI calibration scatter plot};
  \node[dreamlab-grey,font=\sffamily\footnotesize] at (6,2.5) {../output/figures/sci\_calibration\_scatter.png};
\end{tikzpicture}
\caption{SCI predicted condition vs field-surveyed condition (ordinal scale: 1=Poor to 5=Good). Points are jittered for visibility. The diagonal line indicates perfect agreement.}
\label{fig:sci-scatter}
\end{figure}


\section{Model Limitations and Assumptions}

\begin{enumerate}
  \item \textbf{Temporal resolution.} The growing-season median composite integrates spectral information across 5 months. Intra-seasonal management events (mowing, grazing pulses) may not be captured.

  \item \textbf{Spatial resolution.} At 10\,m pixel size, narrow linear features (hedgerows, streams) are represented by mixed pixels. The NDVI transect approach mitigates but does not eliminate this limitation.

  \item \textbf{Training data representativeness.} The PHI-derived training labels may under-represent degraded or transitional habitat states, potentially biasing the classifier towards ``textbook'' class signatures.

  \item \textbf{Condition proxy validity.} The SCI assumes a monotonic relationship between spectral indices and ecological condition. This assumption may break down for habitats where management for biodiversity deliberately reduces biomass (e.g., conservation grazing on species-rich grassland).

  \item \textbf{Independence assumption.} The Monte Carlo simulation treats classification errors as independent between parcels. In practice, spatially correlated classification errors may widen the true confidence interval.
\end{enumerate}

\begin{confidencebox}{MEDIUM}
Despite these limitations, the model provides the best available estimate of biodiversity change given the absence of comprehensive field survey data at both epochs. The identified limitations are systematically addressed through the uncertainty quantification framework.
\end{confidencebox}
